\section{The printed local operator chain and an explicit quotient repair}
\label{sec:source-operator-audit}

This section separates three spaces which the local calculation must
distinguish: the additive Schwartz space on $\mathbb A_{F,S}$,
its $S$-unit coinvariants, and the weighted global Hilbert space
with its orthogonal Paley--Wiener projections.  The first three
assertions below are source-level checks.  The subsequent construction
repairs a particular rational local quotient, with actual operators
and nuclear traces; it does not identify that quotient with the
global weighted projection.

\paragraph{MA-LA-008 (the centering change on pp.~515--516).}
The operator displayed near the bottom of Masullo p.~515 has,
before taking its output coinvariant class, the value
\[
 \begin{aligned}
 A_{\widetilde\varphi}f(x)
  &=\int \widetilde\varphi(g)|g^{-1}x|_S^{1/2}
                     f(g^{-1}x)\,d^\times g
  =|x|_S^{1/2}(Kf)(x),\\
 (Kf)(x)&=\int |g|_S^{-1/2}\widetilde\varphi(g)
                     f(g^{-1}x)\,d^\times g .
 \end{aligned}
 \tag{MA-LA.39}
\]
The second expression, without the output factor $|x|_S^{1/2}$,
is exactly (12.4.1.5) and (12.4.1.7), p.~516.  The paragraph
following (12.4.1.7) calls it the kernel of the same operator.
In coordinates the correct relation is $A=UK$, followed by
the output quotient, where $U=M_{|x|_S^{1/2}}$.
Equivalently $K$ represents $U^{-1}A$ after this explicit
change of target coordinates.  For $F=\mathbb Q$, $S=\{\infty\}$,
and even test functions the finite sign quotient does not remove
the factor: at $x=4$ the two nonzero displayed values differ
by a factor of two.

This is a repairable coordinate omission.  If the centering map
is used to identify operators, all other maps in that diagram
must be conjugated consistently; for example an additive Fourier
cutoff becomes $U^{-1}\mathcal F^{-1}M_b\mathcal F U$ in the
other coordinates.  The two formulas alone do not justify dropping
$U$ while retaining the other unaltered cutoff maps.
Both original pages were inspected as rendered pages, so this
finding is not an extraction ambiguity.

\subsection{A literal failure of the raw coinvariant lift}

Write $S=\{\infty,p\}$, $\mathbb A_S=\mathbb R\times\mathbb Q_p$,
$U_p^0=\mathbb Z_p^\times$, and
\[
 (L_t f)(x)=f(t^{-1}x),\qquad t\in\mathbb Q^\times.
 \tag{MA-LA.40}
\]
For $t=p$ the additive determinant is $|p|_S=1$, so $L_p$
is unitary on additive $L^2(\mathbb A_S)$.  For clarity the
polar-free test space used here is the concrete LF subspace
\[
 \mathcal V=\mathcal S_{\mathrm{flat}}(\mathbb R)
                 \otimes_\iota C_c^\infty(\mathbb Q_p^\times)
       \ \subset\ \mathcal S(\mathbb A_S),
 \qquad
 \mathcal S_{\mathrm{flat}}(\mathbb R)
       =\{u\in\mathcal S(\mathbb R):u^{(j)}(0)=0\ (j\ge0)\}.
 \tag{MA-LA.41}
\]
This spells out the local coordinate-hyperplane conditions used
in Definition~5.3.4.1, p.~259, and its proof, pp.~260--261.
It is stronger than vanishing only at the simultaneous adelic
origin.  Proposition~2.1.4.1(1), pp.~71--72, gives the actual
additive finite tensor structure: every function has a common
finite $p$-adic support and a common additive local-constancy
level.  The additional condition $f(x_\infty,0)=0$ consequently
gives a common neighborhood of $0$ in the $p$-adic variable on
which it vanishes.  Thus each LF stage of (MA-LA.41) has only
finitely many $p$-adic valuations.

There is a notation issue in the source which this explicit
definition avoids.  Remark~2.1.2.1, p.~63, uses
$\mathcal S(\mathbb Q_p^\times)$ for the larger Schwartz space
of the multiplicative group, including infinite rapidly decreasing
valuation sequences.  The proof on pp.~260--261 instead uses the
same notation for compactly supported, origin-vanishing functions
inside additive $\mathcal S(\mathbb Q_p)$, with its finest LF
topology.  No finite-valuation assertion about the larger
multiplicative Schwartz space is used here.

\paragraph{MA-LA-009 (the displayed raw kernel does not kill coinvariant relations).}
Let $q:\mathcal V\to\mathcal V_{\mathbb Z[1/p]^\times}$ be
the separated coinvariant map.  The tensor containment printed
after (12.4.1.7), p.~516, would imply that the raw-output operator
$K:\mathcal V\to\mathcal V$ factors through $q$ on its input.
This implication is false, even for a valid compact trivial-isotype
test supported away from norm one.

\begin{proof}
Choose $0\le h\in C_c^\infty(\mathbb R_{>0})$, nonzero,
supported in an interval disjoint from $1$, and put
\[
 \widetilde\varphi(g_\infty,g_p)
       =\tfrac12 h(|g_\infty|)\mathbf1_{U_p^0}(g_p).
 \tag{MA-LA.42}
\]
For any $g_p\ne0$ exactly one exponent $n$ makes
$p^n g_p$ a unit.  The two signs cancel the factor $1/2$,
so the $S$-unit periodization is exactly
\[
       \sum_{\xi\in\{\pm p^n:n\in\mathbb Z\}}
                \widetilde\varphi(\xi g)
                  =h(|g_\infty||g_p|_p).
 \tag{MA-LA.43}
\]
It is therefore a legitimate lift in (12.4.1.4), p.~515.
Use multiplicative measure with unit mass one at $p$;
other positive Haar normalizations only multiply the next
formula by a nonzero constant.  If $u$ is even and
$f=u\otimes\mathbf1_{U_p^0}$, then
\[
       Kf(x_\infty,x_p)
        =\mathbf1_{U_p^0}(x_p)
             \int_0^\infty h(t)t^{-1/2}u(x_\infty/t)\,\frac{dt}{t}.
 \tag{MA-LA.44}
\]
Take $u\ge0$ to be a nonzero smooth compact annular bump.
The right side is nonzero.  In contrast, $L_pf$ is supported
on the $p$-adic shell $pU_p^0$, and $K L_pf$ is supported on
the same shell, since the lift (MA-LA.42) has unit $p$-adic
component.  Thus $K(f-L_pf)$ is nonzero on $x_p\in U_p^0$,
although $q(f-L_pf)=0$.  The example lies in (MA-LA.41),
and in the source's intended polar-free subspace.
\end{proof}

The valid assertion is that $K$ commutes with the $S$-unit
action and hence induces $qK$ on the quotient.  Equivariance
does not make a raw-output lift constant on input coinvariant
classes.  Taking the output quotient, or specifying a section,
repairs this algebraic point but changes the map whose kernel
and trace have to be justified.

\paragraph{MA-LA-010 (product-norm cutoffs do not preserve additive Schwartz space).}
The Step~II multiplier on p.~516 is not generally an
endomorphism of $\mathcal S(\mathbb A_S)$ when $S$ has a
finite place.  This remains an issue after the proposed
annular position-cutoff repair.

\begin{proof}
Let $b:[0,\infty)\to\mathbb R$ be smooth, equal to one
near zero and zero for large arguments.  Apply the displayed
product-norm multiplier to
\[
       f(x_\infty,x_p)=e^{-x_\infty^2}\mathbf1_{\mathbb Z_p}(x_p).
 \tag{MA-LA.45}
\]
If $b(|x_\infty||x_p|_p)f(x_\infty,x_p)$ belonged to
additive Schwartz space, Proposition~2.1.4.1 would give an
integer $N\ge0$ such that its values at $x_p=0$ and
$x_p=p^N$ agreed for every $x_\infty$.
Choose $x_\infty=r p^N$, with $b(r)=0$.
The first value is $e^{-(r p^N)^2}$ and the second is zero.
This contradicts the required uniform local-constancy level.

For a nonzero smooth annular function $a$, with $a(0)=0$,
choose $r>0$ with $a(r)\ne0$ and use the same test.
The value at $x_p=0$ is zero and the value at $x_p=p^N$
is $a(r)e^{-(r p^N)^2}\ne0$.  Thus the annular multiplier
does not even map all additive Schwartz functions into
additive Schwartz space, and in particular does not have
the printed $\mathcal S\to\mathcal S_{00}$ codomain.
\end{proof}

These functions are perfectly good bounded multipliers on
additive $L^2$.  The failure is their asserted Schwartz-space
domain and codomain, not measurability or boundedness on $L^2$.
It also explains why correcting the real-place inner radius
alone does not complete the mixed-place operator argument.

\subsection{A continuous section and a unitary quotient Fourier transform}

\paragraph{MA-LA-011 (an explicit rational local repair).}
On the actual LF space (MA-LA.41), define
\[
 E_- =\tfrac12(1+L_{-1}),\qquad
 P_0 f=\mathbf1_{U_p^0}(x_p)f,\qquad
 B=P_0\sum_{n\in\mathbb Z}L_{p^n}E_- .
 \tag{MA-LA.46}
\]
This is a continuous projection, its kernel is precisely the
space of $S$-unit coinvariant relations, and $qB=q$.
Its range consists of the functions supported on the unit
valuation shell and invariant under simultaneous negation.
It consequently gives a continuous section of the separated
coinvariant quotient.

\begin{proof}
On a fixed LF stage the valuation support is a finite set.
After $P_0$ only those finitely many $n$ for which $-n$
belongs to that support survive.  Each resulting term is
a continuous dilation and a continuous shell projection.
Thus (MA-LA.46) is continuous on every stage and hence on
the strict LF space.  Reindexing this finite sum shows
$BL_p=B$ and $BL_{-1}=B$.

If $P_j$ denotes the projection onto the shell $p^jU_p^0$,
then, for each $f$, the following sums are finite:
\[
       f=\sum_jP_jf,\qquad
       Bf=\sum_jL_{p^{-j}}P_jE_-f.
 \tag{MA-LA.47}
\]
It follows that $Bf-f$ is a finite sum of unit relations
and hence $qBf=qf$.  On the stated range $B$ is the identity,
so $B^2=B$.  If $Bf=0$, (MA-LA.47) expresses $f$ as a
finite sum of unit relations.  Thus this relation space is
already closed, being $\ker B$, and the induced section is
continuous for the quotient topology.
\end{proof}

The finite-sum assertion must not be reused after additive
Fourier transformation.  With self-dual measures,
\[
 \mathcal F_p\mathbf1_{U_p^0}
       =\mathbf1_{\mathbb Z_p}
                   -p^{-1}\mathbf1_{p^{-1}\mathbb Z_p}.
 \tag{MA-LA.48}
\]
It is nonzero near zero.  Nevertheless its section-periodization
can be calculated and completed on a Hilbert space, rather than
asserted to lie in the original polar-free LF space.

Restrict to the compact trivial isotype, so section vectors are
$u\otimes\mathbf1_{U_p^0}$ with $u$ even.  Let
$\mathcal F_e$ be the usual unitary Fourier transform on
even $L^2(\mathbb R,dx)$ and set
\[
 \begin{aligned}
 (V_pu)(x)&=p^{1/2}u(px),& q_p&=p^{-1/2},\\
 D_p&=(1-q_pV_p^*)(1-q_pV_p)^{-1},&
 H_p&=D_p\mathcal F_e .
 \end{aligned}
 \tag{MA-LA.49}
\]
Then $D_p$ is unitary, $H_p$ is a unitary involution, and
for $u\in\mathcal S_{\mathrm{flat}}(\mathbb R)$ the actual
periodization of $\mathcal F_S(u\otimes\mathbf1_{U_p^0})$
on the unit shell is
\[
 H_pu(x)=(1-p^{-1})\sum_{k\ge0}(\mathcal F_eu)(p^kx)
                         -p^{-1}(\mathcal F_eu)(x/p).
 \tag{MA-LA.50}
\]
The sum converges in $L^2$ and normally with all derivatives
on compact annuli.  It need not be flat at zero.

\begin{proof}
At $x_p\in U_p^0$, (MA-LA.48) gives the coefficient
$1-p^{-1}$ for $n\le0$, the coefficient $-p^{-1}$ for
$n=1$, and zero for $n\ge2$ in the periodization
$\sum_nL_{p^n}\mathcal F_S f$.  This proves the displayed
formula pointwise on annuli.  Rapid decay of $\mathcal F_eu$
proves normal convergence there.  Its $k$th dilation has
$L^2$ norm $p^{-k/2}\|u\|_2$, so the series also converges
absolutely in operator norm.
In terms of $V_p$ its multiplier is
\[
 (1-q_p^2)\sum_{k\ge0}q_p^kV_p^k-q_pV_p^*
       =\frac{1-q_pV_p^*}{1-q_pV_p}=D_p.
 \tag{MA-LA.51}
\]
Since $V_p$ is unitary, numerator and denominator commute
and are adjoints, and the denominator is boundedly invertible.
Thus $D_p^*D_p=1$.
The Fourier scaling identity gives
$\mathcal F_e V_p\mathcal F_e=V_p^*$ and
$\mathcal F_e D_p\mathcal F_e=D_p^{-1}$.
Together with $\mathcal F_e^2=1$, this proves $H_p^2=1$.

For an explicit warning even with zero integral, choose an
even nonnegative annular bump $\psi$ and let
$u(x)=\psi(x)-p^{-1}\psi(x/p)$.
Then $u$ is flat at zero and $\int u=0$.
The series in (MA-LA.50) telescopes to
\[
                H_pu(x)=(\mathcal F_e\psi)(x)
                                -p^{-1}(\mathcal F_e\psi)(x/p),
 \tag{MA-LA.52}
\]
whose value at zero is
$(1-p^{-1})(\mathcal F_e\psi)(0)\ne0$.
Here (MA-LA.52) is the identity on nonzero real annuli,
hence also an $L^2$ identity; its right side gives the
smooth continuation to zero.  Termwise evaluation of
the original periodization at the single point zero
need not equal that continuation.
Thus the intermediate Fourier-periodized function is not
in the flat source space, although it belongs to the stated
Hilbert completion.
\end{proof}

There is now an actual local nuclear construction.  For an even
smooth annular multiplier $a$, a bounded even multiplier $b$,
and $h\in C_c^\infty(\mathbb R_{>0})$, put
\[
 (A_hu)(x)=\int_0^\infty h(t)t^{-1/2}u(x/t)\,\frac{dt}{t},
 \qquad
 T^{\mathrm{quot}}_{a,b,h}=M_a A_h H_pM_bH_p.
 \tag{MA-LA.53}
\]
This is trace class on even $L^2(\mathbb R)$ and factors
nuclearly through the flat unit-shell section.
Indeed the unitary logarithmic map
$u\mapsto\sqrt2e^{y/2}u(e^y)$ converts $M_aA_h$ into
the integral operator with kernel
\[
                  a(e^y)h(e^{y-z}).
 \tag{MA-LA.54}
\]
It is smooth and supported in a compact rectangle.
A Fourier-series expansion on a larger rectangle has coefficients
decreasing faster than every power in both indices.  It gives
an absolutely summable rank-one presentation and, with all output
derivatives, a nuclear map into smooth functions on the compact
output interval.  Multiply each output Fourier vector by a fixed
smooth cutoff supported inside that larger interval and equal to
one on the kernel's output support.  This leaves the kernel
unchanged and makes every output vector extend smoothly by zero.
For the latter nuclear assertion one may choose a
diagonal, faster-than-polynomial scaling of the output Fourier
vectors: the rapid coefficient bounds permit summable functional
norms while the rescaled vectors are bounded in every smooth
seminorm.  Returning from logarithmic coordinates gives a nuclear
map into $\mathcal S_{\mathrm{flat}}(\mathbb R)$.
The remaining factor $H_pM_bH_p$ is bounded, proving the claim.
Insertion into the unit shell and its continuous quotient map
give the trace-compatible factorization of Lemma~C.3.2.2;
the common nuclear rank-one series gives both traces.

This construction is on the quotient Hilbert completion just
specified.  It does not assert that every intermediate operation
preserves additive Schwartz space.  Nor does it assert equality
with the global weighted orthogonal Paley--Wiener projection.
The coupled annular cutoff limit requires its own proof;
MA-LA-014 supplies it below for this rational local sector.
Finite-cutoff nuclearity alone does not fix that limit.

\subsection{The full additive cover cannot supply a nonzero compact trace operator}

\paragraph{MA-LA-012 (unit symmetry prevents compactness on the cover).}
Let $a,b$ be bounded functions of the product norm and let
$K$ be the compactly averaged unitary dilation operator in
(MA-LA.39).  On ordinary additive $L^2(\mathbb A_S)$ the operator
\[
             T=M_{a(|x|_S)}K\mathcal F_S^{-1}
                         M_{b(|\xi|_S)}\mathcal F_S
 \tag{MA-LA.55}
\]
is bounded and commutes with $L_p$.
If it is nonzero it is not compact, hence not trace class.
There are nonzero examples with the valid test (MA-LA.42),
with the printed outer cutoff and also with every nonzero
smooth annular position cutoff.

\begin{proof}
The normalized dilation $|g|_S^{-1/2}f(g^{-1}x)$ is unitary,
so $\|K\|\le\|\widetilde\varphi\|_{L^1(d^\times g)}$.
Both multipliers are bounded and the additive Fourier transform
is unitary.  Norm-dependent multipliers commute with $L_p$,
and multiplicative convolution does also.  Since $|p|_S=1$,
\[
           \mathcal F_S L_p=L_{p^{-1}}\mathcal F_S.
 \tag{MA-LA.56}
\]
This shows that the middle Fourier multiplier commutes with $L_p$.

The powers $L_p^n$ tend weakly to zero on additive $L^2$.
To see this first take two functions with $p$-adic supports
in finite unions of valuation shells; their translated supports
are disjoint for large $n$.  Such functions are dense in $L^2$,
since $x_p=0$ has measure zero, and unitarity gives the general
claim by approximation.  If $T$ were compact, then
$TL_p^nf\to0$ in norm for every $f$.  Commutation instead gives
\[
                     \|TL_p^nf\|_2=\|Tf\|_2.
 \tag{MA-LA.57}
\]
Thus every compact commuting $T$ is zero.

For nonzeroness take $\widetilde\varphi$ as in (MA-LA.42),
and $0\le b\le1$ equal to one near zero and with compact
norm support.  Use the self-dual positive test
$f=e^{-\pi x_\infty^2}\mathbf1_{\mathbb Z_p}(x_p)$.
The function $b\mathcal F_Sf$ is integrable and nonnegative,
with positive integral, so
$v=\mathcal F_S^{-1}(b\mathcal F_S f)$ is continuous and
$v(0)>0$.  Compactness of the dilation parameters and
nonnegativity of the test imply $Kv>0$ in a neighborhood
of the additive origin.  This already proves nonzeroness
for outer position cutoffs, and for sufficiently small
annular inner radii.

For any prescribed nonzero smooth annular $a$, the same
example works.  Gaussian domination shows that $v$, and
then $Kv$, are entire in $x_\infty$ for each fixed $x_p$.
Choose a sufficiently small nonzero $x_p$ so that
$Kv(0,x_p)>0$.  This entire function cannot vanish on an
open real interval.  The set of positive $x_\infty$ for
which $a(|x_\infty||x_p|_p)\ne0$ contains such an interval,
so $M_aKv$ is nonzero.  This completes the example.
\end{proof}

The section-periodization itself has a related continuity
boundary.  If $f_0$ is a nonzero unit-shell vector fixed by
simultaneous negation, then
\[
 f_N=N^{-1/2}\sum_{n=1}^NL_{p^n}f_0,
 \qquad \|f_N\|_2=\|f_0\|_2,
 \qquad Bf_N=\sqrt N f_0.
 \tag{MA-LA.58}
\]
The supports are disjoint valuation shells.  The same argument
works for any norm-dependent additive $L^2$ weight, since
$|p|_S=1$.  Hence the LF section-periodization has no bounded
extension from that full additive Hilbert space.  This does
not contradict the bounded compressed Fourier transform
$H_p$ in MA-LA-011, whose input is the chosen quotient section.

The attribution boundary matters.  The explicit nuclearity
claim on p.~515 concerns a quotient Paley--Wiener level.
On p.~517 Step~III invokes nuclearity after assigning the same
$k^{\mathrm{PW}}$ notation to the raw additive kernel, and
then integrates over a fundamental domain and sums over units.
MA-LA-012 rules out interpreting this as trace class on the
full additive cover.  It does not rule out a quotient trace:
MA-LA-011 constructs one.  A compatible quotient or nuclear
factorization must be supplied before that raw kernel can
be used for the claimed quotient trace.

\subsection{The prime term carried by the repaired quotient transform}

\paragraph{MA-LA-013 (the local all-pass factor has the prime trace).}
In logarithmic coordinates let $\ell=\log p$,
$(V_pf)(x)=f(x+\ell)$, and retain $D_p$ from (MA-LA.49).
Let $g\in C_c^\infty(\mathbb R)$ and let $\phi$ be smooth,
with $\phi(-\infty)=1$, $\phi(+\infty)=0$ and compactly
supported derivative.  Then the actual operator
\[
                  C_g(D_p^*M_\phi D_p-M_\phi)
 \tag{MA-LA.59}
\]
is trace class on $L^2(\mathbb R)$ and
\[
 \operatorname{Tr}\bigl(C_g(D_p^*M_\phi D_p-M_\phi)\bigr)
      = (\log p)\sum_{k\ge1}p^{-k/2}
                    \bigl(g(k\log p)+g(-k\log p)\bigr).
 \tag{MA-LA.60}
\]
For a finite set of primes, $D_S=\prod_{p\in S_f}D_p$
gives the sum of these local prime traces.

\begin{proof}
Write $q=p^{-1/2}$ and the absolutely convergent Laurent series
\[
 D_p=\sum_{n\in\mathbb Z}d_nV_p^n,\qquad
 d_{-1}=-q,\quad d_n=(1-q^2)q^n\ (n\ge0),\quad
 d_n=0\ (n<-1).
 \tag{MA-LA.61}
\]
Unitarity gives $\sum_m\overline d_m d_{m+r}=\delta_{r0}$.
Consequently the difference in (MA-LA.59), before convolution,
is the double sum with coefficients $\overline d_m d_n$ of
\[
    M_{\phi(x-m\ell)-\phi(x)}V_p^{n-m}.
 \tag{MA-LA.62}
\]
Each multiplier difference is compactly supported in an interval
of length $O(1+|m|)$, with all positive-order derivatives
uniformly bounded.  A fixed smooth partition into unit intervals
and the compact smooth kernel of $C_g$ on each piece give
\[
 \bigl\|C_gM_{\phi(x-m\ell)-\phi(x)}V_p^{n-m}\bigr\|_1
                         \le C_{g,\phi,p}(1+|m|).
 \tag{MA-LA.63}
\]
Here translation of a piece is unitary and the final shift
does not change trace norm.  The coefficient series weighted
by $1+|m|$ is summable, proving an actual trace-norm double
series before any traces are taken.

The diagonal kernel of each term gives
$g((n-m)\ell)\int(\phi(x-m\ell)-\phi(x))\,dx$.
The integral equals $m\ell$, by integrating $\phi'$ and
using its total integral $-1$.  Thus the coefficient of
$\ell g(r\ell)$ is
\[
                  a_r=\sum_m m\overline d_m d_{m+r}.
 \tag{MA-LA.64}
\]
For $r\ge1$ its $m=-1$ term is $(1-q^2)q^r$ and its
$m\ge0$ terms sum to $q^{r+2}$, so $a_r=q^r$.
For $r=-k$, $k\ge1$, the $m=k-1$ term is
$-(k-1)(1-q^2)q^k$ and the $m\ge k$ terms sum to
$q^k(q^2+k(1-q^2))$, again giving $a_{-k}=q^k$.
For $r=0$ the two contributions $-q^2$ and $q^2$ cancel.
This proves (MA-LA.60), including both directions and its sign.

All the $D_p$ commute with one another and with $C_g$.
For two factors, expand
\[
 (D_1D_2)^*M_\phi(D_1D_2)-M_\phi
 =D_2^*(D_1^*M_\phi D_1-M_\phi)D_2
                  +(D_2^*M_\phi D_2-M_\phi).
\]
The preceding trace-class estimate and bounded cyclicity make
the trace of the first summand equal to the first local trace.
Induction proves the finite-set assertion.  The same finite
product of (MA-LA.50) computes the compact trivial quotient
Fourier transform for $S=\{\infty\}\cup S_f$.
\end{proof}

Thus the repaired quotient Fourier transform carries the local
prime terms through an explicit trace-class commutator.  This
is an operator statement, not a trace assigned from an explicit
formula.  It does not replace the separate relative-rate analysis
for two compact annular/Fourier cutoffs by a one-sided cutoff.

The same operator calculation has a precise character version.
For $|c|=1$ put
\[
 D_{p,c}=(1-\overline c\,p^{-1/2}V_p^*)
                         (1-cp^{-1/2}V_p)^{-1}.
 \tag{MA-LA.65}
\]
It is unitary.  Its Laurent coefficients are
$d_{-1}=-\overline c\,q$ and
$d_n=(1-q^2)c^nq^n$ for $n\ge0$.
The absolute trace-norm estimates above are unchanged, while
(MA-LA.64) becomes $a_k=c^kq^k$ and
$a_{-k}=\overline c^{\,k}q^k$ for $k\ge1$, with $a_0=0$.
Consequently the trace is
\[
 (\log p)\sum_{k\ge1}p^{-k/2}
       \bigl(c^kg(k\log p)+\overline c^{\,k}g(-k\log p)\bigr).
 \tag{MA-LA.66}
\]
This proves the dual coefficient orientation at the operator
level, in agreement with MA-CC-001 and MA-RS-005.
It is an algebraic generalization of the local unitary factor;
no unproved full character-sector section is inferred from it.

\subsection{The exact remaining source bridge}

Definition~11.2.1.3(2), (11.2.1.8), p.~489, defines
$Q^{\mathrm{PW}}_{S,\Lambda_1,a}$ to be an orthogonal
projection in the weighted global Hilbert space.  Theorem~11.2.1.4
uses it in $P^{\mathrm{geom}}_{\rho,a}A(\varphi)Q^{\mathrm{PW}}$.
The local calculation instead uses a smooth additive Fourier
multiplier.  Literal equality of those two candidates fails
already because a nontrivial smooth multiplier $b$ is not
idempotent, whereas an orthogonal projection is.  This observation
refutes that literal candidate identification; the source does
not actually print such an equality, so it is not attributed to it.

The p.~515 assertion of a linear topological copy also cites
Proposition~5.3.4.4, pp.~263--264.  That proposition states an
equivalence between an entire local zeta transform, existence
of a polar-free representative with the same transform, and
rapid decay of its theta image.  Its displayed statement does
not itself specify the continuous inverse or a nuclear lift
used on p.~515.  The explicit section above supplies such
a map in its stated rational local sector; it is not a claim
that the cited proposition already supplied the full global
operator diagram.

Section~12.3 on p.~514 has no argument, only its heading and
the footnote ``Include section.''  The p.~520 proof then uses
the local theorem to identify the global $k_S$ diagonal integral,
with the two polar terms subtracted.  Proposition~12.2.2.2,
pp.~508--513, is a diagonal integration identity for a
\emph{specified} lifted quotient kernel, under its stated
regularity and absolute-integrability hypotheses.  It supplies
neither the global-to-local map nor the equality of these kernels.
The shell overview on pp.~491--494 also first defines its
valuation cutoffs upstairs on ideles, rather than on idele
classes; a descent or indentation is needed before using them
as quotient kernels.  For example, the condition that a $p$-adic
valuation be zero takes different values on the idele $1$ and
the principal idele $p$, although they represent the same class.

The narrow missing operator statement is therefore a compatible
factorization of the actual condensed global operator through
the local quotient, with the second cyclic composition equal
to the claimed local cutoff operator, together with control
of the limiting remainder and the polar finite-dimensional
part.  Range inclusion alone is insufficient for this use of
Lemma~C.3.2.2.  MA-LA-011 supplies a genuine local nuclear
diagram, and MA-LA-013 supplies its prime-bearing unitary
factor, but neither identifies that diagram with the global
weighted orthogonal projection.  The precise actual global
cutoff value and its lost-divisor criterion are already given
by MA-WP and MA-EG, without adding this missing identification.
The bare inverse obstruction MA-MO and the smoothed
commutators MA-RS retain their separate scopes.

All source locators in this section refer to the inspected
public Masullo draft; no source PDF, text extraction, or
rendered source page is included in the deliverable.
